\begin{foldable} % Knowledge distillation
  The evolution of Artificial Intelligence has resulted in high-capacity models that serve as centralized repositories of specialized expertise.
  However, the substantial size of these models often limits their mobility and deployment in decentralized environments, such as mobile or edge devices.
  To bridge this gap, research on model compression seeks to transfer these expert insights into lightweight, portable agents.
  \textit{Knowledge distillation} (KD)~\cite{hinton2015distilling} is the primary framework for this pedagogical exchange, where an expert network (teacher) guides the learning of a novice network (student).
\end{foldable}

\begin{foldable} % Black-box data-free KD
  In practical AI ecosystems, this information flow is frequently obstructed by systemic barriers that prevent transparent knowledge exchange.
  Stringent privacy regulations, such as HIPAA or GDPR~\cite{us1996hipaa, eu2016gdpr}, prohibit access to original training data, particularly in sensitive domains, such as personal, health, and proprietary data.
  Furthermore, intellectual property or security concerns~\cite{us2016dsta} often render the teacher as \textit{black-box}, withholding internal signals like activations, gradients, or class probabilities, and returning only the top-1 predictions.
  These constraints define the \textit{black-box data-free KD} (BBDFKD) problem: training a student to emulate an expert's intelligence through a narrow interface without any historical data.
\end{foldable}

\begin{foldable} % Recent methods
  Recent methods in BBDFKD~\cite{wang2021zero,zhang2022ideal,yuan2024data} leverage synthetic data, but they still encounter two primary bottlenecks.
  First, they suffer from a lack of domain coverage: synthetic samples fail to replicate the hierarchical structures and semantic variation of natural environments, resulting in a distribution mismatch with the teacher's expertise.
  Second, the distillation signals are suboptimal.
  By relying exclusively on hard top-1 indices, students are denied the \textit{dark} knowledge of inter-class relationships---a cornerstone of KD essential for emulating expert logic.
  Within a collective intelligence framework, this resembles an agent attempting to reconstruct a diverse collective memory from discrete feedback, deprived of the historical data that shaped the expert's knowledge.
\end{foldable}

\begin{foldable} % Our method
  \textbf{Our method.} We propose \textbf{Diverse Image Priors Knowledge Distillation} (\textbf{DIP-KD}) to tackle the BBDFKD problem, \ie, training a student from a black-box teacher with only top-1 signal and without access to any real data.
  To improve KD, we incorporate \textit{diversity}, \ie, how broadly the samples distribute, into a collaborative three-phase pipeline:{
    (1)~\textit{Synthesis}: crafts synthetic \textit{image priors} with diverse patterns and semantics of natural objects;
    (2)~\textit{Contrast}: employs a \textit{primer student} to act as a white-box mediator and ensure synthetic samples are informationally distinct via contrastive learning;
    and (3)~\textit{Distillation}: distills the final student optimally with hard labels from teacher and soft probabilistic labels from the primer student.
  }
\end{foldable}

\begin{foldable} % Contribution
  \textbf{Contribution.} In summary, we highlight our contributions as follows:
  \begin{enumerate}
    \item We introduce \textit{image priors}, a novel class of synthetic images exhibiting diversity across hierarchical structures, nonlinearity, and semantics.
    \item We propose the novel use of a \textit{primer student}, enabling contrastive optimization and extraction of soft probabilistic signals originally restricted in black-box KD.
    \item Extensive analysis of our framework demonstrates that it effectively expands the semantic coverage of synthetic domains, leading to our state-of-the-art performance across 12 benchmarks.
  \end{enumerate}
\end{foldable}
